% Source: https://artofproblemsolving.com/community/c3h2885935_kid_named_lattice_points
% Stable identity: aops-topic-2885935-q01
Near the vertex $(0,0,0)$ we need nonnegative $(x,y,z)$ with $x^2+y^2+z^2=169$. The only ways to write $169$ as a sum of three squares are $13^2+0+0$, $12^2+5^2+0$ and $12^2+4^2+3^2$. Their permutations give $3+6+6=15$ points per vertex, all inside the cube since every coordinate is at most $13$.

No point is counted twice. Being $13$ from two vertices that differ in, say, the $x$-coordinate would need $x\in\{0,3,4,5,12,13\}$ and $25-x\in\{0,3,4,5,12,13\}$, so $x\in\{12,13\}$. With $x=13$, the first vertex forces $y=z=0$, while the second needs a coordinate pair from $\{(5,0),(4,3),\dots\}$ with the $12$ in $x$, which fails; $x=12$ is symmetric. The total is $8\cdot15=\boxed{\textbf{(C)}\ 120}$.
\begin{asy}
usepackage("amsmath");
size(3cm);
defaultpen(fontsize(10pt));
real s=2.5;pair d=(1.1,0.8);pair[] F={(0,0),(s,0),(s,s),(0,s)};draw(F[0]--F[1]--F[2]--F[3]--cycle);draw(F[1]--F[1]+d--F[2]+d--F[3]+d--F[3]);draw(F[2]--F[2]+d);draw(F[0]--F[0]+d--F[1]+d,dashed);draw(F[0]+d--F[3]+d,dashed);draw(arc(F[0],1.3,0,90),red);dot(F[0]);label("$(0,0,0)$",F[0],SW);label("$25$",(s/2,0),S);label("$13$",1.3*dir(45),NE,red);
\end{asy}
